\subsection{Socioeconomic distribution of the change}
\label{subsec:results-socioeconomic}

The aggregate change is small in Curitiba and large in Porto Alegre, and in both cities it is unevenly distributed. Spearman correlations between each axis of the deprivation index and the per-cell increase are positive and significant in every scenario and city (Table~\ref{tab:results-correlation}), which means the cells that lose the most time are also the cells with worse sanitation, lower income and higher illiteracy.

\begin{table}[htbp]
\centering
\caption{Spearman correlation between each axis of the socioeconomic deprivation index and the per-cell increase in access time. All coefficients are significant at $p < 0.001$.}
\label{tab:results-correlation}
\footnotesize
\begin{tabular}{llrrrr}
\hline
City & Scenario & Sanitation & Income & Illiteracy & Composite \\
\hline
Curitiba & \texttt{thermal} & 0.350 & 0.357 & 0.306 & 0.449 \\
Curitiba & \texttt{thermal\_demo} & 0.351 & 0.356 & 0.305 & 0.449 \\
Curitiba & \texttt{thermal+green} & 0.098 & 0.442 & 0.300 & 0.332 \\
Curitiba & \texttt{thermal+green\_demo} & 0.330 & 0.374 & 0.312 & 0.446 \\
Porto Alegre & \texttt{flood} & 0.307 & 0.178 & 0.208 & 0.278 \\
Porto Alegre & \texttt{landslide} & 0.198 & 0.105 & 0.165 & 0.181 \\
Porto Alegre & \texttt{flood+landslide} & 0.530 & 0.336 & 0.398 & 0.504 \\
\hline
\end{tabular}
\end{table}

The composite index and its axes rank the cells differently in each city. In Curitiba the composite carries a higher correlation than any single axis in the two plain thermal scenarios, at 0.449 against 0.357 for income, because the three axes capture different sides of the same condition and their sum separates the affected cells better than each one separately. Adding vegetation breaks that relation. Under \texttt{thermal+green} income alone reaches 0.442 and sanitation falls to 0.098, so the composite drops to 0.332, below the highest axis. The vegetation layer removes the delay from the shaded segments, and the sanitation axis, which described where the heat was mapped, loses its grip on the change in the same movement.

Porto Alegre shows the opposite ordering. Sanitation carries the highest correlation of the three axes in every scenario, at 0.530 under \texttt{flood+landslide}, 0.307 under \texttt{flood} and 0.198 under \texttt{landslide}. The composite stays below it in all three, because income and illiteracy are weakly associated with the change and dilute the signal when given equal weight. The gap widens under the two layers together, where sanitation rises to 0.530 while income falls to 0.336 and illiteracy to 0.398.
